% Presentacion - PFM02 / Sesion 02
% Suma y resta con fracciones

\pdfminorversion=7
\documentclass[
	spanish,
	aspectratio=1610,
	xcolor={dvipsnames,table}
]{beamer}

\geometry{paperwidth=19.2cm,paperheight=12cm}

\usepackage[utf8]{inputenc}
\usepackage[T1]{fontenc}
\usepackage[spanish,es-nodecimaldot]{babel}
\usepackage{lmodern}
\usepackage{booktabs}
\usepackage{array}
\usepackage{graphicx}
\usepackage{ragged2e}
\usepackage{tikz}
\usetikzlibrary{arrows.meta,positioning,calc,decorations.pathreplacing,shapes.geometric,fit,backgrounds,overlay-beamer-styles}
\usepackage{hyperref}

% Datos del alumno y del programa
\newcommand{\studentname}{Mtro. Martin Jonathan de la Cruz Munoz}
\newcommand{\studentshort}{M. de la Cruz Munoz}
\newcommand{\studentid}{00841}
\newcommand{\studentemail}{martin.delacruz@campus.iiiepe.edu.mx}
\newcommand{\programname}{Maestria en Aprendizaje y Ensenanza de las Matematicas}
\newcommand{\programshort}{MAEM}
\newcommand{\coursecode}{PFM02}
\newcommand{\coursename}{Temas selectos de matematicas I}
\newcommand{\coursegroup}{Grupo A}
\newcommand{\courseperiod}{Marzo--julio 2026}
\newcommand{\institutionname}{IIIEPE}
\newcommand{\institutionlongname}{Instituto de Investigacion, Innovacion y Estudios de Posgrado para la Educacion del Estado de Nuevo Leon}
\newcommand{\institutionplace}{Monterrey, Nuevo Leon}
\newcommand{\iiiepelogo}{img/iiiepe/logo-iiiepe-blanco.png}
\newcommand{\iiiepemark}{img/iiiepe/marca-iiiepe.png}
\newcommand{\iiiepedots}{img/iiiepe/logo-puntos.png}

% Paleta institucional
\definecolor{iiiepeNavy}{HTML}{163B5C}
\definecolor{iiiepeBlue}{HTML}{1F6F9F}
\definecolor{iiiepeAqua}{HTML}{20B7A8}
\definecolor{iiiepeOrange}{HTML}{F28C28}
\definecolor{iiiepePaper}{HTML}{F6F8F7}
\definecolor{iiiepeInk}{HTML}{1E2A32}
\definecolor{iiiepeMuted}{HTML}{667085}

\mode<presentation>{
	\usetheme{default}
	\usefonttheme{professionalfonts}
	\setbeamertemplate{navigation symbols}{}
	\setbeamertemplate{blocks}[rounded][shadow=false]
}

\hypersetup{
	colorlinks=true,
	urlcolor=iiiepeBlue,
	linkcolor=iiiepeNavy
}

\setbeamercolor{background canvas}{bg=white}
\setbeamercolor{normal text}{fg=iiiepeInk,bg=white}
\setbeamercolor{structure}{fg=iiiepeBlue}
\setbeamercolor{frametitle}{fg=white,bg=iiiepeNavy}
\setbeamercolor{title}{fg=white,bg=iiiepeNavy}
\setbeamercolor{block title}{fg=white,bg=iiiepeBlue}
\setbeamercolor{block body}{fg=iiiepeInk,bg=iiiepePaper}
\setbeamercolor{alerted text}{fg=iiiepeOrange}
\setbeamerfont{title}{size=\LARGE,series=\bfseries}
\setbeamerfont{frametitle}{size=\large,series=\bfseries}
\setbeamerfont{block title}{series=\bfseries}
\setbeamertemplate{itemize item}{\textcolor{iiiepeAqua}{\large$\blacktriangleright$}}
\setbeamertemplate{itemize subitem}{\textcolor{iiiepeOrange}{\scriptsize$\blacksquare$}}
\setbeamertemplate{enumerate item}{\textcolor{iiiepeBlue}{\insertenumlabel.}}

\newcommand{\localLogo}[1]{%
	\IfFileExists{#1}{\includegraphics[height=0.58cm]{#1}}{\textbf{\institutionname}}%
}

\setbeamertemplate{frametitle}{
	\nointerlineskip
	\begin{beamercolorbox}[wd=\textwidth,ht=1.05cm,dp=0.22cm,leftskip=0.35cm,rightskip=0.25cm]{frametitle}
		\begin{minipage}[c]{0.72\textwidth}
			\insertframetitle
		\end{minipage}
		\hfill
		\raisebox{-0.1cm}{\localLogo{\iiiepelogo}}
	\end{beamercolorbox}
	\vspace{-0.04cm}
	{\color{iiiepeOrange}\rule{\textwidth}{1.4pt}}
}

\setbeamertemplate{footline}{
	\leavevmode
	\hbox{%
		\begin{beamercolorbox}[wd=.34\paperwidth,ht=2.7ex,dp=1ex,leftskip=1em]{author in head/foot}
			\color{white}\studentshort
		\end{beamercolorbox}
		\begin{beamercolorbox}[wd=.38\paperwidth,ht=2.7ex,dp=1ex,center]{title in head/foot}
			\color{white}\programshort\ -- \coursecode
		\end{beamercolorbox}
		\begin{beamercolorbox}[wd=.20\paperwidth,ht=2.7ex,dp=1ex,rightskip=1em plus 1fil]{date in head/foot}
			\color{white}Matricula \studentid\hfill\insertframenumber/\inserttotalframenumber
		\end{beamercolorbox}
	}
}
\setbeamercolor{author in head/foot}{bg=iiiepeNavy}
\setbeamercolor{title in head/foot}{bg=iiiepeBlue}
\setbeamercolor{date in head/foot}{bg=iiiepeNavy}

\newcommand{\accentbar}{%
	\begin{tikzpicture}[remember picture,overlay]
		\path[use as bounding box] (0,0) rectangle (0,0);
		\fill[iiiepeAqua] (current page.south west) rectangle ([xshift=0.42\paperwidth,yshift=0.08cm]current page.south west);
		\fill[iiiepeOrange] ([xshift=0.42\paperwidth]current page.south west) rectangle ([xshift=\paperwidth,yshift=0.08cm]current page.south west);
	\end{tikzpicture}
}

\newcommand{\DividirEntre}[2]{\ensuremath{#1\,\text{ entre }\,#2}}

\tikzset{
	cajaProc/.style={
		draw=iiiepeBlue,
		rounded corners=2mm,
		align=center,
		inner sep=4pt,
		minimum height=8mm,
		font=\small
	},
	cajaDato/.style={
		cajaProc,
		fill=iiiepeBlue!8
	},
	cajaOp/.style={
		cajaProc,
		fill=iiiepeOrange!12
	},
	cajaRes/.style={
		cajaProc,
		fill=iiiepeAqua!12
	},
	flechaProc/.style={
		-{Latex[length=2.4mm]},
		thick,
		draw=iiiepeBlue
	},
	notaProc/.style={
		draw=iiiepeOrange,
		rounded corners=2mm,
		fill=iiiepeOrange!12,
		align=center,
		inner sep=4pt,
		font=\scriptsize
	},
	decisionProc/.style={
		diamond,
		aspect=2.1,
		draw=iiiepeBlue,
		fill=iiiepeAqua!10,
		align=center,
		inner sep=2pt,
		font=\scriptsize
	}
}

\newcommand{\RutaOperacion}[4]{%
\begin{tikzpicture}[node distance=6mm]
	\node[cajaDato,visible on=<1->] (a) {#1};
	\node[cajaOp, right=of a,visible on=<2->] (b) {#2};
	\node[cajaOp, right=of b,visible on=<3->] (c) {#3};
	\node[cajaRes, right=of c,visible on=<4->] (d) {#4};
	\draw[flechaProc,visible on=<1->] (a) -- (b);
	\draw[flechaProc,visible on=<2->] (b) -- (c);
	\draw[flechaProc,visible on=<3->] (c) -- (d);
\end{tikzpicture}%
}

\newcommand{\sectionframe}[1]{%
	\begin{frame}[plain]
		\begin{tikzpicture}[remember picture,overlay]
			\fill[iiiepeNavy] (current page.south west) rectangle (current page.north east);
			\IfFileExists{\iiiepedots}{%
				\node[opacity=0.13,anchor=east] at ([xshift=0.8cm]current page.east) {\includegraphics[height=9.4cm]{\iiiepedots}};
			}{}
			\fill[iiiepeAqua] ([xshift=0.85cm,yshift=2.1cm]current page.south west) rectangle ([xshift=1.05cm,yshift=6.3cm]current page.south west);
			\fill[iiiepeOrange] ([xshift=1.15cm,yshift=2.1cm]current page.south west) rectangle ([xshift=1.35cm,yshift=5.35cm]current page.south west);
		\end{tikzpicture}
		\vspace{1.8cm}
		\hspace{1.25cm}
		\begin{minipage}{0.72\paperwidth}
			{\color{white}\Huge\bfseries #1}\\[0.25cm]
			{\color{white!75}Sesion 02 -- Suma y resta con fracciones}
		\end{minipage}
	\end{frame}
}

% Separadores automaticos desactivados para evitar diapositivas repetitivas.
% \AtBeginSection[]{\sectionframe{\insertsectionhead}}

\title[\coursecode]{Suma y resta con fracciones}
\subtitle{Sesion 02 -- \coursename}
\author[\studentshort]{\texorpdfstring{\studentname\\\footnotesize Matricula: \studentid\\\href{mailto:\studentemail}{\studentemail}}{\studentname, Matricula: \studentid}}
\institute[\institutionname]{\institutionlongname\\\coursecode\ -- \coursegroup}
\date[\courseperiod]{\texorpdfstring{\institutionplace\\\courseperiod}{\institutionplace, \courseperiod}}

\begin{document}

\begin{frame}[plain]
	\begin{tikzpicture}[remember picture,overlay]
		\fill[iiiepeNavy] (current page.south west) rectangle ([xshift=0.56\paperwidth]current page.north west);
		\fill[iiiepeBlue] ([xshift=0.56\paperwidth]current page.south west) rectangle (current page.north east);
		\IfFileExists{\iiiepedots}{%
			\node[opacity=0.16,anchor=south east] at ([xshift=0.8cm,yshift=-0.3cm]current page.south east) {\includegraphics[height=8.6cm]{\iiiepedots}};
		}{}
		\node[anchor=north west] at ([xshift=0.55cm,yshift=-0.45cm]current page.north west) {\localLogo{\iiiepelogo}};
		\fill[iiiepeOrange] ([xshift=0.58\paperwidth,yshift=1.20cm]current page.south west) rectangle ([xshift=0.92\paperwidth,yshift=1.34cm]current page.south west);
		\fill[iiiepeAqua] ([xshift=0.58\paperwidth,yshift=0.90cm]current page.south west) rectangle ([xshift=0.82\paperwidth,yshift=1.04cm]current page.south west);
	\end{tikzpicture}

	\vspace{1.75cm}
	\begin{minipage}{0.51\paperwidth}
		{\color{white}\Huge\bfseries Suma y resta\\ con fracciones}\\[0.25cm]
		{\color{white!86}\large PFM02 -- Sesion 02}\\[0.35cm]
		{\color{iiiepeOrange}\rule{0.82\linewidth}{1.3pt}}\\[0.45cm]
		{\color{white}\studentname}\\
		{\color{white!80}\footnotesize Matricula: \studentid}
	\end{minipage}
	\hfill
	\begin{minipage}{0.35\paperwidth}
		\raggedleft
		{\color{white}\Large\bfseries \institutionname}\\[0.18cm]
		{\color{white!82}\footnotesize \programname}\\[0.5cm]
		{\color{white}\small \coursecode\ -- \coursename}\\
		{\color{white!78}\small \courseperiod}\\[0.5cm]
		{\color{white!74}\footnotesize \institutionplace}
	\end{minipage}
\end{frame}

\begin{frame}{Contenido}
	\tableofcontents
\end{frame}

\section{Actividad}

\begin{frame}{Ficha de la sesion}
	\begin{columns}[T]
		\begin{column}{0.48\textwidth}
			\begin{block}{Contenido}
				Extension del significado de las operaciones y sus relaciones inversas.
			\end{block}
			\begin{block}{PDA}
				Reconoce el significado de las cuatro operaciones basicas y sus relaciones inversas al resolver problemas con numeros con signo.
			\end{block}
		\end{column}
		\begin{column}{0.48\textwidth}
			\begin{block}{Aprendizaje de la secuencia}
				Resolver situaciones problematicas que involucran sumar o restar fracciones.
			\end{block}
			\begin{block}{Producto}
				Procedimientos visibles, ejercicios resueltos, situaciones en contexto, evaluacion y retroalimentacion.
			\end{block}
		\end{column}
	\end{columns}
	\accentbar
\end{frame}

\begin{frame}{Criterios incorporados}
	\small
	\begin{columns}[T]
		\begin{column}{0.49\textwidth}
			\begin{block}{Estructura de secuencia}
				La actividad se organiza en cinco fases: antecedente, explicacion, practica, evaluacion y retroalimentacion.
			\end{block}
			\begin{block}{Evidencia esperada}
				No basta la respuesta: debe observarse estrategia, procedimiento, simplificacion e interpretacion.
			\end{block}
		\end{column}
		\begin{column}{0.49\textwidth}
			\begin{block}{Criterios de evaluacion}
				Actividades diarias, exposicion de secuencia, escrito reflexivo y exposicion final requieren claridad matematica y reflexion.
			\end{block}
			\begin{block}{Decision didactica}
				Cada ejemplo muestra que se hace, por que funciona y como se comprueba.
			\end{block}
		\end{column}
	\end{columns}
	\accentbar
\end{frame}

\begin{frame}{Ruta pedagogica de la actividad}
	\begin{columns}[T]
		\begin{column}{0.20\textwidth}
			\begin{block}{1. Recuperar}
				Simplificar, convertir mixtas e impropias.
			\end{block}
		\end{column}
		\begin{column}{0.20\textwidth}
			\begin{block}{2. Igualar}
				Construir denominador comun.
			\end{block}
		\end{column}
		\begin{column}{0.20\textwidth}
			\begin{block}{3. Operar}
				Sumar o restar numeradores.
			\end{block}
		\end{column}
		\begin{column}{0.20\textwidth}
			\begin{block}{4. Simplificar}
				Reducir a forma irreductible.
			\end{block}
		\end{column}
		\begin{column}{0.20\textwidth}
			\begin{block}{5. Verificar}
				Revisar sentido del contexto.
			\end{block}
		\end{column}
	\end{columns}
	\vspace{0.25cm}
	\begin{block}{Criterio docente}
		El procedimiento es suficiente cuando el estudiante puede explicar la unidad comun, el calculo y la respuesta final.
	\end{block}
	\accentbar
\end{frame}

\begin{frame}{Mapa general de suma y resta con fracciones}
	\centering
	\resizebox{0.98\textwidth}{!}{%
	\begin{tikzpicture}[node distance=0.8cm and 0.95cm]
		\node[cajaDato,visible on=<1->] (inicio) {Fracciones};
		\node[decisionProc, right=of inicio,visible on=<2->] (tipo) {¿hay mixtas\\o enteros?};
		\node[cajaOp, above right=0.45cm and 0.9cm of tipo,visible on=<3->] (convertir) {si:\\convertir a\\impropias};
		\node[cajaOp, below right=0.45cm and 0.9cm of tipo,visible on=<3->] (mantener) {no:\\mantener\\fracciones};
		\node[decisionProc, right=2.65cm of tipo,visible on=<4->] (mismo) {¿mismo\\denominador?};
		\node[cajaOp, above right=0.45cm and 0.9cm of mismo,visible on=<5->] (operar) {operar\\numeradores};
		\node[cajaOp, below right=0.45cm and 0.9cm of mismo,visible on=<5->] (comun) {buscar\\denominador\\comun};
		\node[cajaOp, right=of operar,visible on=<6->] (simplificar) {simplificar};
		\node[cajaRes, right=of simplificar,visible on=<7->] (verificar) {verificar\\contexto};
		\draw[flechaProc,visible on=<1->] (inicio) -- (tipo);
		\draw[flechaProc,visible on=<2->] (tipo) -- (convertir);
		\draw[flechaProc,visible on=<2->] (tipo) -- (mantener);
		\draw[flechaProc,visible on=<3->] (convertir) -- (mismo);
		\draw[flechaProc,visible on=<3->] (mantener) -- (mismo);
		\draw[flechaProc,visible on=<4->] (mismo) -- (operar);
		\draw[flechaProc,visible on=<4->] (mismo) -- (comun);
		\draw[flechaProc,visible on=<5->] (comun.north) -- (operar.south);
		\draw[flechaProc,visible on=<6->] (operar) -- (simplificar);
		\draw[flechaProc,visible on=<7->] (simplificar) -- (verificar);
	\end{tikzpicture}
	}
	\vspace{0.25cm}
	\begin{block}{Lectura del mapa}
		La unidad comun permite comparar partes; la respuesta final debe conservar la cantidad y tener sentido en el problema.
	\end{block}
	\accentbar
\end{frame}

\begin{frame}{Idea clave: partes comparables}
	\begin{columns}[T]
		\begin{column}{0.45\textwidth}
			\begin{block}{Mismo denominador}
				Si las partes tienen el mismo tamano, se operan numeradores y se conserva el denominador.
				\[
					\frac{a}{c}+\frac{b}{c}=\frac{a+b}{c}
				\]
			\end{block}
		\end{column}
		\begin{column}{0.52\textwidth}
			\centering
			\begin{tikzpicture}[scale=0.86]
				\onslide<1->{\foreach \i in {0,1,2,3}{\draw[thick] (\i,0) rectangle +(1,0.45);}}
				\onslide<2->{\fill[iiiepeAqua!65] (0,0) rectangle +(1,0.45);}
				\onslide<3->{\fill[iiiepeAqua!65] (1,0) rectangle +(1,0.45);}
				\onslide<4->{\fill[iiiepeOrange!75] (2,0) rectangle +(1,0.45);}
				\onslide<4->{\node[below] at (2,-0.15) {$\frac{2}{4}+\frac{1}{4}=\frac{3}{4}$};}
				\onslide<4->{\node[above] at (2,0.65) {mismas partes: cuartos};}
			\end{tikzpicture}
		\end{column}
	\end{columns}
	\accentbar
\end{frame}

\section{Procedimientos}

\begin{frame}{Procedimiento 1: mismo denominador}
	\small
	\begin{columns}[T]
		\begin{column}{0.48\textwidth}
			\begin{block}{Suma}
				\[
				\frac{7}{2}+\frac{5}{2}=\frac{7+5}{2}=\frac{12}{2}=6
				\]
				\resizebox{\linewidth}{!}{%
					\RutaOperacion
					{$\frac{7}{2}+\frac{5}{2}$}
					{mismos\\medios}
					{$\frac{7+5}{2}$}
					{$\frac{12}{2}=6$}%
				}
			\end{block}
		\end{column}
		\begin{column}{0.48\textwidth}
			\begin{block}{Resta}
				\[
				\frac{17}{4}-\frac{10}{4}=\frac{17-10}{4}=\frac{7}{4}=1\,\frac{3}{4}
				\]
				\resizebox{\linewidth}{!}{%
					\RutaOperacion
					{$\frac{17}{4}-\frac{10}{4}$}
					{mismos\\cuartos}
					{$\frac{17-10}{4}$}
					{$\frac{7}{4}=1\,\frac{3}{4}$}%
				}
			\end{block}
		\end{column}
	\end{columns}
	\begin{block}{Verificacion}
		El denominador no se suma ni se resta porque representa el tamano de la parte.
	\end{block}
	\accentbar
\end{frame}

\begin{frame}{Procedimiento 2: denominador comun}
	\small
	\begin{columns}[T]
		\begin{column}{0.46\textwidth}
			\begin{block}{Paso 1: elegir unidad comun}
				Para $\frac{3}{4}+\frac{1}{2}$, el mcm de $4$ y $2$ es $4$.
			\end{block}
			\begin{block}{Paso 2: convertir}
				\[
				\frac{1}{2}=\frac{2}{4}
				\]
			\end{block}
		\end{column}
		\begin{column}{0.50\textwidth}
			\begin{block}{Paso 3: operar}
				\[
				\frac{3}{4}+\frac{2}{4}=\frac{5}{4}=1\,\frac{1}{4}
				\]
			\end{block}
			\begin{block}{Sentido}
				$\frac{3}{4}$ mas $\frac{1}{2}$ debe ser mayor que $1$; $1\,\frac{1}{4}$ cumple esa estimacion.
			\end{block}
		\end{column}
	\end{columns}
	\vspace{0.15cm}
	\centering
	\resizebox{0.86\textwidth}{!}{%
		\RutaOperacion
		{$\frac{3}{4}+\frac{1}{2}$}
		{igualar\\denominador}
		{$\frac{3}{4}+\frac{2}{4}$}
		{$\frac{5}{4}=1\,\frac{1}{4}$}%
	}
	\vspace{0.12cm}
	\begin{block}{Nota}
		El denominador comun no cambia la cantidad; solo expresa ambas fracciones con la misma unidad.
	\end{block}
	\accentbar
\end{frame}

\begin{frame}{Procedimiento 3: producto cruzado}
	\small
	\begin{columns}[T]
		\begin{column}{0.48\textwidth}
			\begin{block}{Regla como atajo}
				\[
				\frac{a}{b}+\frac{c}{d}=\frac{ad+bc}{bd}
				\]
				\[
				\frac{a}{b}-\frac{c}{d}=\frac{ad-bc}{bd}
				\]
			\end{block}
		\end{column}
		\begin{column}{0.48\textwidth}
			\begin{block}{Ejemplo}
				\[
				\frac{1}{2}+\frac{3}{7}=\frac{1(7)+3(2)}{14}=\frac{13}{14}
				\]
			\end{block}
			\begin{block}{Cuidado}
				Despues del producto cruzado siempre se revisa si la fraccion puede simplificarse.
			\end{block}
		\end{column}
	\end{columns}
	\vspace{0.15cm}
	\centering
	\resizebox{0.78\textwidth}{!}{%
		\RutaOperacion
		{$\frac{1}{2}+\frac{3}{7}$}
		{crear comun\\$2\cdot7=14$}
		{$\frac{1(7)+3(2)}{14}$}
		{$\frac{13}{14}$}%
	}
	\vspace{0.12cm}
	\begin{block}{Sentido del atajo}
		El producto cruzado crea un denominador comun usando el producto de los denominadores.
	\end{block}
	\accentbar
\end{frame}

\begin{frame}{Procedimiento 4: enteros y fracciones mixtas}
	\small
	\begin{columns}[T]
		\begin{column}{0.48\textwidth}
			\begin{block}{Convertir mixta}
				\[
				5\,\frac{2}{3}=\frac{5(3)+2}{3}=\frac{17}{3}
				\]
			\end{block}
			\begin{block}{Convertir entero}
				\[
				4=\frac{4}{1}
				\]
			\end{block}
		\end{column}
		\begin{column}{0.48\textwidth}
			\begin{block}{Ejemplo completo}
				\[
				5\,\frac{2}{3}+\frac{2}{8}=\frac{17}{3}+\frac{1}{4}
				\]
				\[
				\frac{68}{12}+\frac{3}{12}=\frac{71}{12}=5\,\frac{11}{12}
				\]
			\end{block}
		\end{column}
	\end{columns}
	\vspace{0.12cm}
	\centering
	\resizebox{0.72\textwidth}{!}{%
		\RutaOperacion
		{$5\,\frac{2}{3}$}
		{multiplicar\\entero}
		{$\frac{5(3)+2}{3}$}
		{$\frac{17}{3}$}%
	}
	\vspace{0.08cm}
	\resizebox{0.82\textwidth}{!}{%
		\RutaOperacion
		{$\frac{17}{3}+\frac{1}{4}$}
		{denominador\\$12$}
		{$\frac{68}{12}+\frac{3}{12}$}
		{$\frac{71}{12}=5\,\frac{11}{12}$}%
	}
	\accentbar
\end{frame}

\section{Ejercicios resueltos}

\begin{frame}{Antecedente: simplificar fracciones}
	\scriptsize
	La tabla muestra varios incisos resueltos; el esquema representa la ruta comun que se aplica al bloque.
	\vspace{0.08cm}
	\begin{tabular}{llll}
		\toprule
		\textbf{Fraccion} & \textbf{Resultado} & \textbf{Fraccion} & \textbf{Resultado}\\
		\midrule
		$\frac{4}{20}$ & $\frac{4/4}{20/4}=\frac15$ & $\frac{9}{27}$ & $\frac{9/9}{27/9}=\frac13$\\[0.15cm]
		$\frac{6}{24}$ & $\frac{6/6}{24/6}=\frac14$ & $\frac{24}{32}$ & $\frac{24/8}{32/8}=\frac34$\\[0.15cm]
		$\frac{12}{30}$ & $\frac{12/6}{30/6}=\frac25$ & $\frac{15}{25}$ & $\frac{15/5}{25/5}=\frac35$\\
		\bottomrule
	\end{tabular}
	\vspace{0.25cm}
	\centering
	\resizebox{0.62\textwidth}{!}{%
		\RutaOperacion
		{$\frac{4}{20}$}
		{dividir ambos\\entre $4$}
		{$\frac{4/4}{20/4}$}
		{$\frac{1}{5}$}%
	}
	\vspace{0.14cm}
	\begin{block}{Criterio}
		Una fraccion esta simplificada cuando numerador y denominador ya no comparten divisor mayor que $1$.
	\end{block}
	\accentbar
\end{frame}

\begin{frame}{Antecedente: impropias y mixtas}
	\scriptsize
	La tabla muestra varios incisos resueltos; los esquemas representan la ruta comun que se aplica al bloque.
	\vspace{0.08cm}
	\begin{columns}[T]
		\begin{column}{0.49\textwidth}
			\begin{block}{Impropias a mixtas}
				\begin{tabular}{ll}
				$\frac54=1\,\frac14$ & $\frac{50}{3}=16\,\frac23$\\
				$\frac83=2\,\frac23$ & $\frac{11}{2}=5\,\frac12$\\
				$\frac74=1\,\frac34$ & $\frac43=1\,\frac13$\\
				\end{tabular}
			\end{block}
		\end{column}
		\begin{column}{0.49\textwidth}
			\begin{block}{Mixtas a impropias}
				\begin{tabular}{ll}
				$5\,\frac32=\frac{13}{2}$ & $15\,\frac13=\frac{46}{3}$\\
				$9\,\frac23=\frac{29}{3}$ & $7\,\frac78=\frac{63}{8}$\\
				$8\,\frac34=\frac{35}{4}$ & $2\,\frac{3}{16}=\frac{35}{16}$\\
				\end{tabular}
			\end{block}
		\end{column}
	\end{columns}
	\vspace{0.1cm}
	\centering
	\resizebox{0.58\textwidth}{!}{%
		\RutaOperacion
		{$\frac{50}{3}$}
		{$50=16(3)+2$}
		{$16+\frac{2}{3}$}
		{$16\,\frac{2}{3}$}%
	}
	\quad
	\resizebox{0.36\textwidth}{!}{%
		\RutaOperacion
		{$15\,\frac{1}{3}$}
		{$15(3)+1$}
		{$\frac{46}{3}$}
		{$\frac{46}{3}$}%
	}
	\accentbar
\end{frame}

\begin{frame}{Ejercicios: mismo denominador}
	\small
	La tabla muestra varios incisos resueltos; el esquema representa la ruta comun que se aplica al bloque.
	\vspace{0.08cm}
	\begin{block}{Operar numeradores y conservar denominador}
	\[
	\frac{7}{2}+\frac{5}{2}=\frac{12}{2}=6
	\qquad
	\frac{8}{7}+\frac{1}{7}+\frac{2}{7}=\frac{11}{7}=1\,\frac47
	\]
	\[
	\frac{17}{4}-\frac{10}{4}=\frac{7}{4}=1\,\frac34
	\qquad
	\frac{12}{9}-\frac{3}{9}-\frac{4}{9}=\frac{5}{9}
	\]
	\end{block}
	\centering
	\resizebox{0.72\textwidth}{!}{%
		\RutaOperacion
		{$\frac{8}{7}+\frac{1}{7}+\frac{2}{7}$}
		{mismos\\septimos}
		{$\frac{8+1+2}{7}$}
		{$\frac{11}{7}=1\,\frac{4}{7}$}%
	}
	\vspace{0.1cm}
	\begin{block}{Revision}
		Todas las respuestas deben estar simplificadas o expresadas como mixta cuando convenga para interpretar la cantidad.
	\end{block}
	\accentbar
\end{frame}

\begin{frame}{Ejercicios: distinto denominador por producto cruzado}
	\small
	La tabla muestra varios incisos resueltos; el esquema representa la ruta comun que se aplica al bloque.
	\vspace{0.08cm}
	\begin{block}{Producto cruzado como denominador comun}
	\[
	\frac12+\frac37=\frac{7+6}{14}=\frac{13}{14}
	\qquad
	\frac69+\frac{2}{10}=\frac{60+18}{90}=\frac{13}{15}
	\]
	\[
	\frac32-\frac35=\frac{15-6}{10}=\frac{9}{10}
	\qquad
	\frac{10}{3}-\frac94=\frac{40-27}{12}=\frac{13}{12}=1\,\frac{1}{12}
	\]
	\end{block}
	\centering
	\resizebox{0.78\textwidth}{!}{%
		\RutaOperacion
		{$\frac{1}{2}+\frac{3}{7}$}
		{cruzar\\productos}
		{$\frac{1(7)+3(2)}{14}$}
		{$\frac{13}{14}$}%
	}
	\accentbar
\end{frame}

\begin{frame}{Ejercicios: distinto denominador por equivalentes}
	\small
	La tabla muestra varios incisos resueltos; el esquema representa la ruta comun que se aplica al bloque.
	\vspace{0.08cm}
	\begin{block}{Construir fracciones equivalentes}
	\[
	\frac34+\frac12=\frac54=1\,\frac14
	\qquad
	\frac74+\frac{3}{20}=\frac{38}{20}=\frac{19}{10}
	\]
	\[
	\frac23+\frac17=\frac{17}{21}
	\qquad
	\frac49-\frac13=\frac19
	\]
	\[
	\frac{7}{15}-\frac25=\frac{1}{15}
	\qquad
	\frac53-\frac{4}{27}=\frac{41}{27}=1\,\frac{14}{27}
	\]
	\end{block}
	\centering
	\resizebox{0.78\textwidth}{!}{%
		\RutaOperacion
		{$\frac{3}{4}+\frac{1}{2}$}
		{convertir\\$\frac{1}{2}=\frac{2}{4}$}
		{$\frac{3}{4}+\frac{2}{4}$}
		{$\frac{5}{4}=1\,\frac{1}{4}$}%
	}
	\accentbar
\end{frame}

\section{Situaciones}

\begin{frame}{Situaciones en contexto: combinar cantidades}
	\small
	\begin{columns}[T]
		\begin{column}{0.48\textwidth}
			\begin{block}{01 Receta}
				Escalar de $4$ a $12$ porciones: $1\,\frac12\cdot3=4\,\frac12$ tazas de agua; $\frac14\cdot3=\frac34$ de piloncillo.
			\end{block}
			\begin{block}{02 Pizza}
				\[
				\frac13+\frac14=\frac{4}{12}+\frac{3}{12}=\frac{7}{12}
				\]
			\end{block}
		\end{column}
		\begin{column}{0.48\textwidth}
			\begin{block}{04 Ciclismo}
				\[
				\frac28+5\,\frac23=\frac14+\frac{17}{3}=\frac{71}{12}=5\,\frac{11}{12}
				\]
			\end{block}
			\begin{block}{08 Distancia}
				\[
				4\,\frac12+3\,\frac58+3\,\frac34=\frac{95}{8}=11\,\frac78
				\]
			\end{block}
		\end{column}
	\end{columns}
	\accentbar
\end{frame}

\begin{frame}{Situaciones en contexto: encontrar lo que falta}
	\small
	\begin{columns}[T]
		\begin{column}{0.48\textwidth}
			\begin{block}{03 Caramelos}
				\[
				1-\bigl(\frac14+\frac12\bigr)=1-\frac34=\frac14
				\]
			\end{block}
			\begin{block}{05 Yeso}
				\[
				1\,\frac12-\frac64=\frac32-\frac32=0
				\]
			\end{block}
		\end{column}
		\begin{column}{0.48\textwidth}
			\begin{block}{06 Rancho}
				\[
				1-\bigl(\frac18+\frac14\bigr)=1-\frac38=\frac58
				\]
			\end{block}
			\begin{block}{07 Encuesta}
				\[
				1-\bigl(\frac14+\frac16+\frac13\bigr)=\frac14
				\]
			\end{block}
		\end{column}
	\end{columns}
	\accentbar
\end{frame}

\begin{frame}{Situacion 09: tablas apiladas}
	\small
	\begin{columns}[T]
		\begin{column}{0.46\textwidth}
			\begin{block}{Sumar espesores}
				\[
				\frac{7}{12}+\frac56+\frac23
				\]
				\[
				\frac{7}{12}+\frac{10}{12}+\frac{8}{12}=\frac{25}{12}=2\,\frac{1}{12}
				\]
			\end{block}
			\resizebox{\linewidth}{!}{%
				\RutaOperacion
				{$\frac{7}{12}+\frac56+\frac23$}
				{convertir\\a doceavos}
				{$\frac{7+10+8}{12}$}
				{$2\,\frac{1}{12}$}%
			}
		\end{column}
		\begin{column}{0.50\textwidth}
			\centering
			\begin{tikzpicture}[scale=0.82]
				\foreach \y/\h/\lab/\col in {0/0.45/$\frac{2}{3}$/iiiepeBlue!35,0.62/0.56/$\frac{5}{6}$/iiiepeOrange!45,1.35/0.39/$\frac{7}{12}$/iiiepeAqua!45}{
					\draw[thick, fill=\col] (0,\y) rectangle (4,\y+\h);
					\draw[thick, dashed] (-0.45,\y) -- (0,\y);
					\draw[thick, dashed] (-0.45,\y+\h) -- (0,\y+\h);
					\node[left] at (-0.5,\y+0.5*\h) {\lab};
				}
				\node[notaProc] at (2,2.25) {sumar espesores};
				\draw[flechaProc] (2,2.02) -- (2,1.80);
			\end{tikzpicture}
		\end{column}
	\end{columns}
	\accentbar
\end{frame}

\begin{frame}{Situacion 10: perimetro de rectangulo}
	\begin{columns}[T]
		\begin{column}{0.45\textwidth}
			\begin{block}{Datos}
				Largo $4\,\frac34=\frac{19}{4}$ cm; ancho $2\,\frac12=\frac{5}{2}$ cm.
			\end{block}
			\begin{block}{Procedimiento}
				\[
				P=2\biggl(\frac{19}{4}+\frac{5}{2}\biggr)
				\]
				\[
				P=2\biggl(\frac{19}{4}+\frac{10}{4}\biggr)
				=2\biggl(\frac{29}{4}\biggr)=\frac{29}{2}
				\]
			\end{block}
			\resizebox{\linewidth}{!}{%
				\RutaOperacion
				{$l+a$}
				{igualar\\cuartos}
				{$\frac{19}{4}+\frac{10}{4}$}
				{$\frac{29}{4}$}%
			}
		\end{column}
		\begin{column}{0.52\textwidth}
			\centering
			\begin{tikzpicture}[scale=0.86]
				\draw[thick, fill=iiiepeAqua!12] (0,0) rectangle (5.2,2.4);
				\draw[thick, <->] (0,-0.35) -- (5.2,-0.35);
				\draw[thick, <->] (5.55,0) -- (5.55,2.4);
				\node[below] at (2.6,-0.35) {$4\,\frac{3}{4}$ cm};
				\node[right] at (5.55,1.2) {$2\,\frac{1}{2}$ cm};
				\node at (2.6,1.2) {$P=\frac{29}{2}$ cm};
			\end{tikzpicture}
			\vspace{0.2cm}
			\begin{block}{Respuesta}
				Opcion a: $\frac{29}{2}$.
			\end{block}
		\end{column}
	\end{columns}
	\accentbar
\end{frame}

\section{Evaluacion}

\begin{frame}{Autoevaluacion y evidencia}
	\scriptsize
	\begin{tabular}{p{0.58\textwidth}ccc}
		\toprule
		\textbf{Indicador} & \textbf{3} & \textbf{2} & \textbf{1}\\
		\midrule
		Realizo sumas de fracciones usando el algoritmo aprendido. & $\square$ & $\square$ & $\square$\\[0.16cm]
		Realizo restas de fracciones de forma correcta. & $\square$ & $\square$ & $\square$\\[0.16cm]
		Identifico la operacion necesaria para resolver un problema. & $\square$ & $\square$ & $\square$\\[0.16cm]
		Justifico denominador comun, simplificacion y respuesta contextual. & $\square$ & $\square$ & $\square$\\
		\bottomrule
	\end{tabular}
	\vspace{0.25cm}
	\begin{block}{Cierre metacognitivo}
		Que error debo vigilar: sumar denominadores, no simplificar, operar sin convertir mixtas o no interpretar el contexto.
	\end{block}
	\accentbar
\end{frame}

\begin{frame}{Retroalimentacion}
	\small
	\begin{columns}[T]
		\begin{column}{0.49\textwidth}
			\begin{block}{Errores frecuentes}
				\begin{itemize}
					\item Sumar denominadores.
					\item Omitir la simplificacion.
					\item Tratar mixtas como si fueran enteros separados.
					\item Responder sin unidad o contexto.
				\end{itemize}
			\end{block}
		\end{column}
		\begin{column}{0.49\textwidth}
			\begin{block}{Preguntas de revision}
				\begin{itemize}
					\item Que unidad comun use?
					\item El resultado puede simplificarse?
					\item La respuesta es razonable?
					\item La operacion corresponde al problema?
				\end{itemize}
			\end{block}
		\end{column}
	\end{columns}
	\accentbar
\end{frame}

\section{Cierre}

\begin{frame}{Conclusiones}
	\begin{enumerate}
		\item Sumar o restar fracciones exige operar partes comparables.
		\item El denominador comun representa una unidad compartida, no un adorno del algoritmo.
		\item Las fracciones mixtas se convierten a impropias para evitar errores de operacion.
		\item La respuesta debe simplificarse e interpretarse en el contexto.
		\item La actividad sigue las cinco fases institucionales: antecedente, explicacion, practica, evaluacion y retroalimentacion.
	\end{enumerate}
	\accentbar
\end{frame}

\begin{frame}[plain]
	\begin{tikzpicture}[remember picture,overlay]
		\fill[iiiepeNavy] (current page.south west) rectangle (current page.north east);
		\IfFileExists{\iiiepedots}{%
			\node[opacity=0.13,anchor=south east] at ([xshift=0.8cm,yshift=-0.2cm]current page.south east) {\includegraphics[height=9cm]{\iiiepedots}};
		}{}
		\node[anchor=north west] at ([xshift=0.55cm,yshift=-0.45cm]current page.north west) {\localLogo{\iiiepelogo}};
	\end{tikzpicture}
	\vspace{2.2cm}
	\begin{center}
		{\color{white}\Huge\bfseries Gracias}\\[0.55cm]
		{\color{white}\Large \studentname}\\[0.18cm]
		{\color{white!78}\footnotesize Matricula \studentid\ -- \studentemail}\\[0.45cm]
		{\color{iiiepeOrange}\rule{0.36\paperwidth}{1.4pt}}\\[0.35cm]
		{\color{white!72}\footnotesize \institutionname\ -- \coursecode\ -- \coursename}
	\end{center}
\end{frame}

\end{document}

